\section{Robusthet}
\label{sec:robust}

En parser som bara fungerar när allt går rätt är inte en parser. Tre fall måste hanteras, och alla
tre inträffar på riktigt.

\subsection{Skräpdata före SOF}

Noden kopplas in mitt i en pågående överföring, eller en störning lägger till bytes. Strömmen ser
ut så här:

\begin{console}
32 74 00 FF A5 F7 03 03 25 17 7F 05 10 20 30 02 C2
|---------| |-------------------------------------|
  skräp                    framen
\end{console}

Lösningen faller ut ur tillståndsmaskinen gratis: i \code{WaitForSof1} kastas varje byte som inte
är \code{0xA5}. Parsern står stilla tills strömmen börjar se ut som en frame.

Ett fall är värt att tänka igenom särskilt: parsern står i \code{WaitForSof2} och får en byte som
inte är \code{0xF7}. Det naiva svaret är att gå tillbaka till \code{WaitForSof1} och kasta byten.
Men om byten var \code{0xA5} har man just kastat starten på en riktig frame. Rätt svar är att
stanna i \code{WaitForSof2}: den byte man nyss såg kan vara den nya frames första.

\subsection{Tappade bytes}

En byte försvinner. Då blir alla följande fält förskjutna ett steg, längdfältet läser en databyte
som längd, och parsern väntar på fel antal bytes.

Det finns ingen elegant räddning, och det ska inte finnas någon: framen \emph{är} förlorad.
Poängen är att parsern ska upptäcka det och komma tillbaka till ett känt läge, inte att den ska
gissa. Checksumman fallerar, \code{extractFrame()} returnerar \code{false}, anroparen anropar
\code{reset()}, och nästa \code{0xA5F7} i strömmen startar om.

\subsection{Orimlig längd}

Det här är fallet som är lätt att missa, och det enda i kapitlet som är en säkerhetsbugg och inte
bara en felhantering.

Längdfältet är en byte och kan alltså vara upp till 255. Parserns buffert rymmer
\code{MaxPayloadLen} bytes. Tar parsern emot en längd som är större än så, och sedan lydigt
skriver in så många databytes, skriver den utanför sin egen buffert.

\begin{cppcode}
case State::WaitForLen:
{
    // Reject a length the buffer cannot hold, else the payload would overflow myBuf.
    if (MaxPayloadLen < byte) { return reportParseError(); }

    myBuf[myParsedBytes++] = byte;
    myState                = State::WaitForType;
    break;
}
\end{cppcode}

\begin{aside}
\note[Regel] Varje längd som kommer utifrån ska kontrolleras mot buffertens storlek \emph{innan}
den används, inte efter. Regeln gäller överallt i den här boken, och den gäller
\code{TX_DLC} i \kapref{ch:driver} av exakt samma skäl.
\end{aside}

\subsection{Två frames efter varandra}

När en frame är klar och extraherad måste parsern nollställas innan nästa kan tas emot. Glöms det
står den kvar i \code{Ready} och kastar varje ny byte.

Det är därför \code{reset()} finns i det publika gränssnittet, och det är därför noden i
\kapref{ch:bussar} anropar den efter varje extraherad frame.
